\section{Descripción general}

\subsection{Perspectiva del producto}
Resilis es una aplicación web, accesible desde el celular, que acompaña la ayuda humanitaria desde que una persona la pide hasta que la recibe. Funciona en dos niveles. En el nivel de \emph{sectores} decide dónde concentrar los esfuerzos. En el nivel de \emph{solicitudes} identifica a cada persona que necesita algo, la prioriza, verifica su situación, asigna la entrega y confirma que la ayuda llegó, como en una aplicación de pedidos: el solicitante ve el estado de su solicitud y recibe un código que debe mostrar al equipo de entrega para recibir la ayuda. Del lado de la oferta, registra las donaciones y el inventario, y calcula cómo repartir los recursos disponibles entre las solicitudes para cubrir la mayor necesidad posible con lo que hay.

Antes del desastre el sistema guarda una \emph{línea base} por sector (viviendas, población, vulnerabilidad), construida con datos proyectados y con datos oficiales abiertos del Estado. Después del desastre recibe reportes de daños y solicitudes de ayuda, y recalcula la prioridad y la confianza. Los datos de personas, daños y necesidades del prototipo son simulados, anclados en cifras reales.

\subsection{Usuarios y roles}
Todas las personas que reportan o piden ayuda tienen un usuario registrado, porque el sistema debe poder dar seguimiento y verificar. El peso de un reporte depende de quién lo envía: una persona natural tiene un peso bajo y un brigadista un peso alto, porque su cargo le exige reportar.

\begin{table}[H]
\centering
\caption{Roles del sistema y peso de la fuente.}
\begin{tabularx}{\textwidth}{@{}p{2.8cm}Xl@{}}
\toprule
\textbf{Rol} & \textbf{Descripción y permisos principales} & \textbf{Peso de fuente $F$} \\
\midrule
Solicitante (ciudadano) & Persona registrada de la comunidad. Envía reportes de daños y solicitudes de ayuda, ve solo sus propios datos, el estado de su solicitud y su código de entrega. & 0,40 \\
Brigadista & Personal de campo con obligación de reportar. Envía y verifica reportes y solicitudes, crea solicitudes en nombre de otras personas y puede actuar como equipo de entrega. & 0,90 \\
Equipo de entrega & Brigadistas y organizaciones registradas. Ven solo las entregas que les asignaron, validan el código y registran la entrega. & 0,90 brigadista; 0,70 organización \\
Donante & Persona, organización o empresa registrada. Registra donaciones, puede dirigirlas, y ve el estado y el destino agregado de cada una, sin datos personales de los beneficiarios. & no reporta \\
Centro de acopio & Organización registrada que recibe donaciones, verifica su cantidad y estado, y gestiona el inventario. & 0,70 \\
Entidad coordinadora &Organismos encargados de organizar la respuesta. Ven todo, ajustan pesos, aprueban el plan de asignación de recursos y asignan entregas. & 0,95 \\
Administrador & Gestiona usuarios, roles, datos base y parámetros. No ve los datos de vulnerabilidad de las personas. & no reporta \\
\bottomrule
\end{tabularx}
\end{table}

Los \emph{datos oficiales del Estado} (reportes de eventos y de daños) no son un usuario, sino una fuente con peso propio, propuesto en 0,85. Los valores de $F$ son una propuesta inicial y se pueden ajustar (sección~\ref{sec:confianza}).

\subsection{Entorno de operación}
Los usuarios de campo usan celulares con conexión móvil limitada o intermitente. Las entregas deben poder confirmarse \textbf{sin conexión}: el equipo de entrega valida el código en su dispositivo y sincroniza después. Quien no tiene celular recibe una ficha impresa con su código, y si no hay electricidad el reporte o la solicitud se hace en papel y un brigadista la digitaliza. Los usuarios de coordinación trabajan en un computador o tableta. El sistema se despliega en la nube en Railway.

\subsection{Restricciones}
\begin{itemize}
  \item \textbf{Tiempo:} 12 horas de desarrollo, a cargo de una sola persona.
  \item \textbf{Costo:} sin gasto adicional para el prototipo. Se usa software abierto y la cuenta de pago de Railway que el equipo ya tiene.
  \item \textbf{Tecnología:} Python como lenguaje principal; aplicación web; despliegue en Railway.
  \item \textbf{Datos:} los daños, las necesidades y las personas del prototipo son simulados. Los datos oficiales se cargan desde archivos abiertos.
  \item \textbf{Legales:} protección de datos personales según la Ley 1581 de 2012. Los datos sobre discapacidad o salud son sensibles y requieren autorización expresa.
  \item \textbf{Concurso:} todo el material entregado pasa a dominio público.
\end{itemize}

\subsection{Supuestos y dependencias}
\begin{itemize}
  \item Las proyecciones de viviendas por municipio del Atlántico sirven como base para generar sectores simulados. Aún se están recopilando los datos del departamento.
  \item Existen datos abiertos oficiales de eventos y daños que se pueden cargar como archivos. Las fuentes concretas están por identificar y confirmar.
  \item Se dispone de mapas abiertos (OpenStreetMap) para la visualización.
  \item El servicio de Railway está disponible durante la demostración.
\end{itemize}

\subsection{Alternativas descartadas}
Durante el diseño se analizaron y descartaron estas alternativas:
\begin{itemize}
  \item \textbf{Depender de satélites, drones, redes sociales y sensores.} Quedan como mejora futura y opcional, porque el reto parte de comunicaciones caídas.
  \item \textbf{Limitarse a un solo desastre y zona} (inundación en Barranquilla, o vendavales en Malambo y Soledad). Se amplió a los desastres naturales del Atlántico, porque lo único que cambia entre eventos son los pesos y los escenarios.
  \item \textbf{Un modelo opaco} que entregue solo un resultado. Se eligió un índice multicriterio explicable.
  \item \textbf{Bajar la prioridad de los sectores de difícil acceso.} Se tratan como alerta para no dejar atrás a las comunidades más aisladas.
  \item \textbf{Reportes anónimos.} Se descartaron porque no permiten dar seguimiento ni verificar. Todo reporte y toda solicitud se asocian a un usuario registrado.
\end{itemize}

\subsection{Fuera del alcance de esta versión}\label{sec:fuera}
No se incluyen: conexión automática (por interfaz de programación) con sistemas oficiales como el de la UNGRD, imágenes satelitales o de drones, análisis de redes sociales, sensores, modelos supervisados de predicción, mensajería por SMS o WhatsApp, el procesamiento de pagos de donaciones monetarias (solo se registran), verificación de identidad contra bases oficiales y uso operativo real. La carga de archivos abiertos de datos oficiales sí está incluida. Los demás elementos se mencionan como línea de crecimiento, porque el reto parte de comunicaciones caídas y el sistema no debe depender de ellos.
